\documentclass[12pt]{article}
\usepackage[utf8]{inputenc}
\usepackage[T1]{fontenc}
\usepackage[margin=1in]{geometry}
\usepackage{setspace}
\usepackage{xcolor}
\usepackage[colorlinks=true, urlcolor=blue!60!black, linkcolor=blue!60!black]{hyperref}
\setlength{\parindent}{0em}
\setlength{\parskip}{0.6em}
\renewcommand{\baselinestretch}{1.0}
\begin{document}
\begin{center}
{\LARGE \textit{Central Hudson Gas \& Electric Corp. v. Public Service Commission}, 447 U.S. 557 (1980)\par}
\vspace{0.5em}
{\large Case Brief --- Two-Page Edition\par}
\end{center}

\section{Memory Jogger}
New York's ban on utility promotional advertising fell 8--1 under a new four-part test for commercial speech (\href{https://supreme.justia.com/cases/federal/us/447/557/}{Justia}).

\section{Detailed Case Facts}
During the 1973 oil embargo, the New York Public Service Commission banned electric utilities' advertising ``promot[ing] the use of electricity.'' In 1977 it extended the ban on promotional ads indefinitely, citing energy conservation and rate inequities, while allowing informational ads. Central Hudson Gas \& Electric Corp., a regulated monopoly, challenged the ban.

\section{Procedural History}
The New York courts upheld the order, reasoning that a monopoly's ads in a ``noncompetitive market'' had little value. The Supreme Court reversed in 1980.

\section{Judicial Votes}
8--1. Powell wrote for Burger, Stewart, White, and Marshall; Brennan, Blackmun, and Stevens concurred in the judgment; Rehnquist dissented.

\section{Holding}
Commercial speech receives lesser protection, but restrictions on truthful speech about lawful activity must serve a substantial interest, directly advance it, and be no more extensive than necessary. The total ban failed the last requirement.

\section{Analysis of Opinions}
\textbf{Majority:} The ads were truthful and lawful, and utilities face interfuel competition. Conservation was a substantial interest with an ``immediate connection'' to advertising, though the rate rationale was ``highly speculative.'' But the ban reached ads for efficient devices such as heat pumps, and narrower measures were available. \textbf{Blackmun:} suppressing truthful information to ``dampen'' demand ``strikes at the heart of the First Amendment.'' \textbf{Stevens:} the definition of commercial speech is too broad. \textbf{Rehnquist:} the decision revives \textit{Lochner}.

\section{Examples: Future Applications}
\textbf{Same-Side:} a ban on advertising liquor prices fails, as in \textit{44 Liquormart}. \textbf{Opposite-Side:} ads for betting that is illegal in the State are unprotected under the first prong. \textbf{Fence-Sitter:} a billboard ban for marijuana that is legal under state law but illegal under federal law.

\section{Critique}
Powell offered two inconsistent definitions of commercial speech, a boundary problem that persists. Following Rehnquist, critics warn that aggressive review lets courts deregulate consumer protection and public health under the banner of free speech.

\section{Key Quotations}
\begin{quote}
``In commercial speech cases, then, a four-part analysis has developed. At the outset, we must determine whether the expression is protected by the First Amendment. For commercial speech to come within that provision, it at least must concern lawful activity and not be misleading. Next, we ask whether the asserted governmental interest is substantial. If both inquiries yield positive answers, we must determine whether the regulation directly advances the governmental interest asserted, and whether it is not more extensive than is necessary to serve that interest.''
\end{quote}
\noindent --- Justice Powell, for the Court, 447 U.S. at 566.

\end{document}
